\documentclass[11pt]{article}
\usepackage[a4paper,margin=26mm]{geometry}
\usepackage[T1]{fontenc}
\usepackage{lmodern,microtype}
\usepackage{amsmath,amssymb,amsthm,mathtools,bm}
\usepackage{booktabs,array,longtable}
\usepackage{xurl}
\usepackage[hidelinks,pdftitle={Sharper bounds for RE-03},pdfsubject={Informally audited partial research results; the full query complexity is not determined}]{hyperref}
\usepackage{enumitem,needspace,listings}
\setlength{\parindent}{0pt}
\setlength{\parskip}{5.5pt}
\setlength{\emergencystretch}{2em}
\setcounter{tocdepth}{1}
\newtheorem{theorem}{Theorem}[section]
\newtheorem{lemma}[theorem]{Lemma}
\newtheorem{proposition}[theorem]{Proposition}
\newtheorem{corollary}[theorem]{Corollary}
\theoremstyle{remark}
\newtheorem{remark}[theorem]{Remark}
\newcommand{\R}{\mathbb R}
\newcommand{\E}{\mathbb E}
\newcommand{\Prob}{\mathbb P}
\newcommand{\F}{\mathsf F}
\newcommand{\op}{\mathrm{op}}
\newcommand{\Hclass}{\mathcal H_{n,k}}
\newcommand{\OPT}{\operatorname{OPT}}
\newcommand{\rank}{\operatorname{rank}}
\newcommand{\tr}{\operatorname{tr}}
\newcommand{\KL}{D_{\mathrm{KL}}}
\newcommand{\dist}{\operatorname{dist}}
\newcommand{\eps}{\varepsilon}
\newcommand{\norm}[1]{\left\lVert #1\right\rVert}
\newcommand{\ip}[2]{\left\langle #1,#2\right\rangle}
\newcommand{\Id}{I}
\DeclareMathOperator{\diag}{diag}
\DeclareMathOperator{\spanop}{span}
\DeclareMathOperator{\orth}{orth}
\DeclareMathOperator{\col}{col}
\lstset{basicstyle=\ttfamily\small,columns=fullflexible,keepspaces=true,
 breaklines=true,frame=none,showstringspaces=false}
\title{Sharper bounds for RE-03}
\author{Sidney Holden\\{\small Center for Computational Biology, Flatiron Institute}\\{\small Simons Foundation}}
\date{}
\begin{document}
\maketitle
\vspace{-1.6em}
\begin{abstract}
For the fixed-partition HODLR approximation problem with $n=2^Lk$, this note
develops the bounds
\[
 c\min\left\{n,\frac{kL}{\eps}+\frac{k}{\eps^2}\right\}
 \le q_*(n,k,\eps)\le
 C\min\left\{n,\frac{kL^2}{\eps}+\frac{kL}{\eps^2}\right\}.
\]
The proposed upper bound is obtained by a row-isotropic perturbation inequality
for truncated SVD, a noisy Gaussian Nystr\"om calculation, a constant-accuracy perforated pilot, and a globally shared final regression. The proposed new lower term $kL/\eps$ uses weighted
hierarchical Stiefel frames, large fixed rank-constrained anchors, and a
transcript-divergence argument valid for adaptive queries on both sides.
The earlier $k/\eps^2$ lower-bound argument is included in an appendix, so the
combined argument can be read without an external unpublished dependency.
These bounds still do not match in general. They retain the characterization
$q_* = \Theta(n)$ for $\eps\le\sqrt{k/n}$ and rule out the earlier candidate
rate $kL+k/\eps^2$ as a uniform answer.
\end{abstract}

\textbf{Status (13 September 2026).} This is a partial research manuscript,
not a full solution of RE-03. A separate Codex AI agent passed an independent
informal audit of the stated partial proofs; see the accompanying
\texttt{independent-review.md}. No formal verification or external human peer
review is asserted. The numerical checks are sanity checks, not proof
certificates; the original continuation suite is incomplete in the supplied
archive, and only the available assembly smoke checks were rerun.
No matching algorithm for the lower expression is supplied.

\clearpage
\begingroup\small\setlength{\parskip}{0pt}\tableofcontents\endgroup
\clearpage

\section{Problem, principal bounds, and scope}

Let $n=2^Lk$, where $k\ge1$ and $L\ge2$ are integers, and let
$0<\eps<1/2$. The class $\Hclass$ consists of matrices whose two ordered
off-diagonal sibling blocks at each node of the prescribed balanced binary
partition have rank at most $k$. The terminal diagonal blocks of order $k$
are unrestricted. Set
\[
 \OPT(A)=\min_{H\in\Hclass}\norm{A-H}_{\F}.
\]
One query returns $Av$ or $A^{\mathsf T}v$. The side and vector may depend
measurably on earlier answers and independent randomness. Only vector queries
are charged; a matrix with $s$ query columns costs $s$ calls. The output must
be in $\Hclass$ and satisfy
\begin{equation}\label{eq:guarantee}
 \Prob\{\norm{A-B}_{\F}\le(1+\eps)\OPT(A)\}\ge0.99
 \quad\text{for every fixed }A\in\R^{n\times n}.
\end{equation}
This is the quantifier order and exact-real arithmetic model of RE-03.\footnote{
RE-03, \emph{Optimal matvec query complexity of HODLR approximation}, problem
statement and September 10, 2026 audit; source~\cite{re03}.}

The repository cites an upper bound $O(\min\{n,kL^4/\eps^3\})$ and a lower
bound $\Omega(kL+k/\eps)$ in the regime $n\ge c_0k/\eps$.\footnote{
Chen et al., \emph{Near-optimal hierarchical matrix approximation from
matrix-vector products}, Theorems 1.1--1.2 and Section 6; source~\cite{chen}.}
The earlier supplied note~\cite{prior} proposed the stronger accuracy lower
bound $c_*\min\{n,k/\eps^2\}$ and an exact-recovery obstruction $k(L+1)$.
Its relevant arguments are preserved in Appendix~\ref{app:accuracy} and
Section~\ref{sec:geometry}.

\begin{theorem}[Combined bounds]\label{thm:combined}
For every admissible $(n,k,\eps)$, the arguments in this manuscript give
\begin{equation}\label{eq:combined}
 c\min\left\{n,\frac{kL}{\eps}+\frac{k}{\eps^2}\right\}
 \le q_*(n,k,\eps)
 \le C\min\left\{n,\frac{kL^2}{\eps}+\frac{kL}{\eps^2}\right\},
\end{equation}
where $c,C>0$ are universal. One conservative lower constant is
\[
 D_0=2^{24},\qquad c_*=(6400D_0^2)^{-1},\qquad c=c_*/8.
\]
The upper bound is nonadaptive and admits the explicit query budget in
Theorem~\ref{thm:twostage}. Independently, $q_*\ge k(L+1)$.
\end{theorem}

The constants are deliberately loose. An asymptotic bound with these constants
is not a claim that the reference implementation is efficient at moderate
matrix sizes. In particular, its theorem-parameter mode normally selects the
$n$-query fallback on small numerical examples.

\begin{corollary}[Full-recovery regime]\label{cor:full}
For $0<\eps\le\sqrt{k/n}$,
\[
 c_* n\le q_*(n,k,\eps)\le n.
\]
Thus $q_*=\Theta(n)$ in this regime with universal constants.
\end{corollary}

\begin{corollary}[A previously plausible smaller rate is excluded]
There is no universal bound
$q_*=O(\min\{n,kL+k/\eps^2\})$.
\end{corollary}
\begin{proof}
Take $k=1$, $n=2^L$, and $\eps=L^{-1/2}$ for $L>4$.
The lower expression in \eqref{eq:combined} is at least a constant times
$L^{3/2}$, while $L+\eps^{-2}=2L$. Both are below $n$ for large $L$.
Their ratio is unbounded.
\end{proof}

The remaining gap is substantive. At fixed positive accuracy and increasing
$L$, the expressions in \eqref{eq:combined} scale as $kL$ and $kL^2$ before
capping. At $\eps=L^{-1/2}$ they scale as $kL^{3/2}$ and $kL^{5/2}$.
Neither the $n$-query fallback nor the saturation corollary closes these gaps.

\section{Block geometry and the exact baseline}\label{sec:geometry}

Let $\mathcal R$ be the collection of all ordered sibling blocks $(I,J)$ and
$\mathcal D$ the terminal diagonal blocks. These sets partition the entries.
For a rectangular matrix $C$, let $[C]_k$ be a fixed measurable choice of a
best rank-at-most-$k$ SVD truncation and write
$\tau_k(C)^2=\sum_{i>k}\sigma_i(C)^2$.

\begin{lemma}[Separable optimum]\label{lem:separable}
For every matrix $A$,
\begin{equation}\label{eq:separable}
 \OPT(A)^2=\sum_{(I,J)\in\mathcal R}\tau_k(A_{I,J})^2.
\end{equation}
For any $B\in\Hclass$, each quantity
$\norm{A_{I,J}-B_{I,J}}_\F^2-\tau_k(A_{I,J})^2$ is nonnegative,
as are the diagonal-block errors. Their sum is
$\norm{A-B}_\F^2-\OPT(A)^2$.
\end{lemma}
\begin{proof}
The squared Frobenius norm is the sum of the squared norms on disjoint blocks.
The rank constraints on different sibling blocks are independent. Minimize
on each block by an SVD, and copy the unrestricted diagonal blocks.
\end{proof}

\begin{proposition}[The exact $n$-query upper bound]\label{prop:upper}
Querying $Ae_j$ for $j=1,\ldots,n$ and then applying the blockwise projection
of Lemma~\ref{lem:separable} returns an exactly optimal matrix in $\Hclass$.
\end{proposition}

\begin{proposition}[Adaptive-safe exact-recovery obstruction]\label{prop:exactlower}
Every algorithm satisfying \eqref{eq:guarantee} uses at least $k(L+1)$ queries
in the worst case.
\end{proposition}
\begin{proof}
In each ordered sibling block allow only its first $k$ columns to be nonzero,
and allow arbitrary entries in the $k\times k$ diagonal leaves. This is a
linear subspace $\mathcal V\subseteq\Hclass$. At each of the $L$ levels its
sibling blocks contribute $nk$ free entries, and the leaves contribute another
$nk$, so $\dim\mathcal V=nk(L+1)$.

Place an independent standard Gaussian prior on coordinates in an orthonormal
basis of $\mathcal V$. Conditional on the algorithm's seed and previous
answers, a new vector query supplies at most $n$ linear scalar observations.
By induction, the unobserved component remains a nondegenerate Gaussian on a
subspace of dimension at least $nk(L+1)-nq$. Adaptivity does not add
constraints beyond those answers: the next linear measurement is fixed once
the past is fixed. If $q<k(L+1)$, the posterior has no atoms, and any measurable
point estimate has conditional probability zero of exact recovery.

But $\OPT(A)=0$ on this prior, so \eqref{eq:guarantee} requires exact recovery
with probability at least $0.99$ for each fixed input and hence on average.
This contradicts the preceding posterior statement.
\end{proof}

\section{A perturbation bound for rank truncation}\label{sec:truncation}

The following deterministic inequality is the main ingredient of the improved
upper bound. It separates the effect of noise in the leading left singular
space from the full Frobenius size of the noise.

\begin{lemma}[Deterministic truncation inequality]\label{lem:truncation}
Let $X,E\in\R^{r\times m}$ and $1\le k\le\min(r,m)$. Let $P$ be a rank-$k$
projector onto leading left singular vectors of $X$, set $T=(I-P)X$ and
$\tau=\norm{T}_\F$. Then
\begin{align}\label{eq:truncation}
 \norm{X-[X+E]_k}_\F^2-\tau^2
 &\le 4\tau\norm{PE}_\F+6\norm{E}_\F^2
       +4\sqrt{2k}\norm{T^{\mathsf T}E}_\F.
\end{align}
No independence, spectral gap, or distributional hypothesis is needed.
\end{lemma}
\begin{proof}
Put $F=PE$, $G=(I-P)E$, $Z=X+F$, and $f=\norm{F}_\F$.
Let $Q$ and $Q_0$ be rank-$k$ right singular projectors for $Z+G$ and $Z$,
respectively. The optimality of $Q$ gives
\[
 \norm{Z(I-Q)}_\F^2\le \tau_k(Z)^2+a,
 \qquad
 a=2\sqrt{2k}\norm{T^{\mathsf T}G}_\F+\norm{G}_\F^2.
\]
Indeed, expand $(Z+G)^{\mathsf T}(Z+G)$ in the variational characterization of
its leading right singular subspace. The cross term is bounded using
$Z^{\mathsf T}G=T^{\mathsf T}G$ and
$\norm{Q-Q_0}_\F\le\sqrt{2k}$. The remaining term obeys
$\tr(G^{\mathsf T}G(Q-Q_0))\le\norm{G}_\F^2$.

Since $\tau_k(Z)\le\tau+f$,
\begin{align*}
 \norm{X(I-Q)}_\F^2
 &\le\bigl(\sqrt{(\tau+f)^2+a}+f\bigr)^2\\
 &\le\tau^2+4\tau f+5f^2+2a.
\end{align*}
The second line follows from $\sqrt{x^2+a}\le x+\sqrt a$ and
$2f\sqrt a\le f^2+a$. Also $[X+E]_k=(X+E)Q$, and right-orthogonality gives
\[
 \norm{X-(X+E)Q}_\F^2
 =\norm{X(I-Q)}_\F^2+\norm{EQ}_\F^2.
\]
Finally $T^{\mathsf T}G=T^{\mathsf T}E$ and
$5\norm{F}_\F^2+2\norm{G}_\F^2+\norm{E}_\F^2\le6\norm{E}_\F^2$.
\end{proof}

\begin{corollary}[Row-isotropic noise]\label{cor:isotropic}
Suppose the second moments of $E$ satisfy
\[
 \E[E_{i\alpha}E_{j\beta}]=h\,\delta_{ij}\Sigma_{\alpha\beta},
 \qquad \Sigma\succeq0,
\]
and put $S=\tr\Sigma$. Then
\begin{equation}\label{eq:isotropic}
 \E\norm{X-[X+E]_k}_\F^2
 \le \tau_k(X)^2+10\tau_k(X)\sqrt{khS}+6rhS.
\end{equation}
\end{corollary}
\begin{proof}
The three second moments needed in Lemma~\ref{lem:truncation} are
\[
 \E\norm{PE}_\F^2=khS,\quad
 \E\norm{T^{\mathsf T}E}_\F^2=h\norm{T}_\F^2S,\quad
 \E\norm{E}_\F^2=rhS.
\]
Apply Cauchy--Schwarz to the two unsquared norms and use
$4+4\sqrt2<10$. If $r<k$, truncate at $r$ instead: the risk is simply
$rhS$, and the upper expression with $\tau_k(X)=0$ remains valid.
\end{proof}

\begin{lemma}[Gaussian least-squares moments]\label{lem:regression}
Let $G\in\R^{t\times r}$ be standard Gaussian, $t>r+1$, and let
$F\in\R^{t\times m}$ be independent of $G$ with independent centered Gaussian
rows of covariance $\Sigma$. Then $E=G^\dagger F$ has second moments
\[
 \E[E_{i\alpha}E_{j\beta}]
 =\frac{\delta_{ij}\Sigma_{\alpha\beta}}{t-r-1}.
\]
\end{lemma}
\begin{proof}
Conditional on $G$, the row covariance factor is $(G^{\mathsf T}G)^{-1}$.
Sign changes of individual columns of $G$ make the expected off-diagonal
entries of this inverse zero. A Schur complement expresses each diagonal
entry as the reciprocal squared norm of one Gaussian column projected onto
the orthogonal complement of the others. That squared norm is
$\chi^2_{t-r+1}$. Integrating its density gives
$\E[(\chi^2_\nu)^{-1}]=1/(\nu-2)$ for $\nu>2$.
Thus $\E(G^{\mathsf T}G)^{-1}=I/(t-r-1)$, proving the claim.
\end{proof}

\section{Noisy Gaussian Nystr\"om approximation}\label{sec:nystrom}

Consider a fixed block $C\in\R^{m_1\times m_2}$ with additive sketch
contaminants $M\in\R^{m_1\times a}$ and $N\in\R^{b\times m_2}$.
Draw independent standard Gaussian matrices of the following sizes:
\[
 \Omega:m_2\times s,\quad \Gamma:a\times s,\quad
 \Psi:m_1\times t,\quad \Theta:b\times t.
\]
Form
\begin{equation}\label{eq:gn}
 Q=\orth(C\Omega+M\Gamma),\quad
 \widehat X=(\Psi^{\mathsf T}Q)^\dagger
              (\Psi^{\mathsf T}C+\Theta^{\mathsf T}N),\quad
 \widehat C=Q[\widehat X]_k.
\end{equation}
The number of columns of $Q$ may be less than $s$.

\begin{theorem}[Noisy block estimate]\label{thm:gn}
Assume $s>k+1$ and $t>s+1$. Set
\[
 \tau=\tau_k(C),\quad \alpha=\frac{k}{s-k-1},\quad d=t-s-1,
 \quad D=(1+\alpha)\tau^2+\alpha\norm M_\F^2+\norm N_\F^2.
\]
Then the construction \eqref{eq:gn} satisfies
\begin{align}\label{eq:gnbound}
 \E\norm{C-\widehat C}_\F^2
 \le{}&\tau^2+\alpha(\tau^2+\norm M_\F^2)\\
 &+10\tau\sqrt{\frac{kD}{d}}+\frac{6sD}{d}.\nonumber
\end{align}
The estimate is rank at most $k$.
\end{theorem}
\begin{proof}
If $k\ge\min(m_1,m_2)$, the same arguments below can use a smaller true rank
in the right-space calculation; alternatively use its rank-$k$ SVD with
zero singular directions completed orthogonally. We give the case
$k\le\min(m_1,m_2)$ explicitly.

Write $C=C_k+C_{\rm tail}$ using an SVD and let $V_1$ contain its leading
$k$ right singular vectors. Let $\Omega_1=V_1^{\mathsf T}\Omega$ and complete
$V_1$ to an orthogonal basis. Because $\Omega_1$ has full row rank almost
surely, multiplication of the range sketch by $\Omega_1^\dagger$ gives
\[
 \norm{(I-QQ^{\mathsf T})C_k}_\F^2
 \le\norm{(C_{\rm tail}\Omega+M\Gamma)\Omega_1^\dagger}_\F^2.
\]
The Gaussian coordinates entering $C_{\rm tail}\Omega$ and $\Gamma$ are
independent of $\Omega_1$. Their cross term has expectation zero. The
Gaussian pseudoinverse calculation of Lemma~\ref{lem:regression}, transposed,
(a standard ingredient in Gaussian range finding~\cite{halko}),
therefore gives
\begin{equation}\label{eq:rangemiss}
 \E\norm{(I-QQ^{\mathsf T})C_k}_\F^2
 \le\alpha(\tau^2+\norm M_\F^2).
\end{equation}
One may verify the omitted right isometry in the first display by multiplying
on the right by $V_1^{\mathsf T}$; it does not change the Frobenius norm.

The ideal projected rank-$k$ approximation satisfies
\begin{align}\label{eq:idealrange}
 \norm{C-Q[Q^{\mathsf T}C]_k}_\F^2
 &\le\norm{C-QQ^{\mathsf T}C_k}_\F^2\\
 &=\tau^2+\norm{(I-QQ^{\mathsf T})C_k}_\F^2.\nonumber
\end{align}
The equality follows from orthogonality of the right singular spaces of
$C_k$ and $C_{\rm tail}$. In particular, for $C_\perp=(I-QQ^{\mathsf T})C$,
\[
 \E\norm{C_\perp}_\F^2
 \le(1+\alpha)\tau^2+\alpha\norm M_\F^2.
\]

Condition on $Q$ and let $r$ be its number of columns. The Gaussian matrix
$\Psi^{\mathsf T}Q$ is independent of $\Psi^{\mathsf T}C_\perp$.
Consequently
\[
 \widehat X=Q^{\mathsf T}C+E,
 \quad
 E=(\Psi^{\mathsf T}Q)^\dagger
      (\Psi^{\mathsf T}C_\perp+\Theta^{\mathsf T}N)
\]
has the second moments of Lemma~\ref{lem:regression} with
$\Sigma=C_\perp^{\mathsf T}C_\perp+N^{\mathsf T}N$ and
$h=1/(t-r-1)$. Apply Corollary~\ref{cor:isotropic} with
$X=Q^{\mathsf T}C$. Its tail is at most $\tau$, because projection is a
contraction and $Q^{\mathsf T}C_k$ is an admissible rank-$k$ competitor.
The error outside $Q$ is orthogonal to the error inside $Q$.
Use \eqref{eq:idealrange}, $r\le s$, $t-r-1\ge d$, and Jensen's inequality
on the square root after averaging over $Q$. This gives \eqref{eq:gnbound}.
If $r=0$, the output is zero and the same bound follows from
\eqref{eq:idealrange}; no inverse is needed.
\end{proof}

\section{Perforated peeling with a global error bound}\label{sec:algorithm}

Only the first $h=L-1$ levels need nontrivial rank truncation. The remaining
size-$2k$ diagonal blocks are unrestricted in $\Hclass$, since their
$k\times k$ sibling blocks already have rank at most $k$.

\subsection{Sketches and query accounting}

At a level $\ell\in\{1,\ldots,h\}$, the target blocks have order
$m=n/2^\ell$. Treat the two orientations separately. For the upper
orientation let $(I_b,J_b)$ be the left/right children of parent $b$.
Let $H_{<\ell}$ contain the approximations of the strictly coarser sibling
blocks, with zero entries elsewhere, and write $R=A-H_{<\ell}$.
Draw a uniform independent hash $f(b)\in\{1,\ldots,B\}$ and independent
Gaussian $m\times s$ matrices $\Omega_b$. For each bucket $a$, query the
matrix whose restriction to $J_b$ is $\Omega_b$ when $f(b)=a$, and is zero
elsewhere. Subtract the known product with $H_{<\ell}$.
The rows $I_b$ of the bucket $f(b)$ give
\begin{equation}\label{eq:rightsketch}
 Y_b=A_{I_b,J_b}\Omega_b+
      \sum_{c\ne b:f(c)=f(b)}R_{I_b,J_c}\Omega_c.
\end{equation}
This uses at most $Bs$ vector queries for one orientation.

Independently draw a left hash $g$ and Gaussian $m\times t$ matrices
$\Psi_b$. Query $A^{\mathsf T}$ on bucketed vectors supported on the $I_b$
and subtract $H_{<\ell}^{\mathsf T}$. Transpose the appropriate outputs to get
\begin{equation}\label{eq:leftsketch}
 Z_b=\Psi_b^{\mathsf T}A_{I_b,J_b}+
      \sum_{c\ne b:g(c)=g(b)}\Psi_c^{\mathsf T}R_{I_c,J_b}.
\end{equation}
This costs at most $Bt$ queries. Return
$Q_b[(\Psi_b^{\mathsf T}Q_b)^\dagger Z_b]_k$ on this block, where
$Q_b=\orth(Y_b)$. The lower orientation exchanges the two children and uses
fresh independent randomness. Freeze $H_{<\ell}$ until both orientations
are complete.

Every queried vector depends only on random sketches, hashes, $n$, and $k$;
subtraction of earlier approximations occurs after the oracle responses.
Thus all actual query vectors may be chosen in advance. Empty buckets may
be skipped. The per-level cost is at most $2B(s+t)$.

\subsection{Contamination energies}

Let
\[
 a_\ell=\sum_{\text{ordered blocks at level }\ell}\tau_k(A_{I,J})^2,
 \qquad O=\sum_{\ell=1}^{h}a_\ell=\OPT(A)^2.
\]
Let $T_{\ell-1}$ be the squared error on the strictly coarser blocks.
Condition on the past. For a right sketch the contaminant $M_b$ is the
horizontal concatenation of the coarser residual blocks that collide with
$b$; for a left sketch $N_b$ is their vertical analogue. Concatenated
Gaussian matrices put \eqref{eq:rightsketch}--\eqref{eq:leftsketch} exactly
in the setting of Theorem~\ref{thm:gn}.

Every pair of distinct parents collides with probability $1/B$. Moreover,
the corresponding residual entries belong to strictly coarser blocks.
Across both orientations each such entry is counted at most once in the
right-contamination sum, and at most once in the left-contamination sum.
Therefore, with the expectation here only over hashes,
\begin{equation}\label{eq:contamination}
 \E\sum_b\norm{M_b}_\F^2\le T_{\ell-1}/B,
 \qquad
 \E\sum_b\norm{N_b}_\F^2\le T_{\ell-1}/B.
\end{equation}
The two sums need not be independent. No independence between different
output blocks is needed in the analysis.

\subsection{Summing the errors before choosing parameters}

Put
\[
 \alpha=\frac{k}{s-k-1},\quad
 \beta=\frac{s}{t-s-1},\quad
 \gamma=\sqrt{\frac{k}{t-s-1}},
\]
\[
 r_0=\alpha+6\beta(1+\alpha),\qquad
 c_0=10\gamma\sqrt{1+\alpha}.
\]
Summing Theorem~\ref{thm:gn} over all targets at this level, using
\eqref{eq:contamination}, and applying Cauchy--Schwarz to the mixed terms
gives
\begin{align}\label{eq:levelbound}
 \E[\text{level-}\ell\text{ error}^2\mid\text{past}]
 \le{}&(1+r_0+c_0)a_\ell+
       \frac{r_0}{B}T_{\ell-1}\\
 &+\frac{c_0}{\sqrt B}\sqrt{a_\ell T_{\ell-1}}.\nonumber
\end{align}
For clarity, the total quantity $D$ in the mixed term of the block theorem
has conditional expectation at most
$(1+\alpha)(a_\ell+T_{\ell-1}/B)$.
Use
$\sqrt{a(a+x)}\le a+\sqrt{ax}$.
The coefficient on $T_{\ell-1}/B$ from the non-mixed terms is no greater
than $r_0$, giving \eqref{eq:levelbound}.

Choose $B=h$ and let $S$ be the expected squared error on all nonterminal
sibling blocks. Since $T_{\ell-1}$ is a sum of earlier nonnegative block
errors, $\sum_\ell\E T_{\ell-1}\le hS$. Also
\[
 \sum_\ell\E\sqrt{a_\ell T_{\ell-1}}
 \le\sqrt{O\sum_\ell\E T_{\ell-1}}
 \le\sqrt{OhS}.
\]
Consequently
\begin{equation}\label{eq:globalineq}
 S\le(1+r_0+c_0)O+r_0S+c_0\sqrt{OS}.
\end{equation}
All expectations are finite: Theorem~\ref{thm:gn} and induction on the finite
number of levels establish this before the summation argument.

\subsection{Parameters and the last diagonal fit}

Set
\begin{equation}\label{eq:parameters}
 \eta=\eps/2600,\quad
 s=k+1+\lceil k/\eta\rceil,\quad
 t=s+1+\lceil s/\eta\rceil,\quad B=h.
\end{equation}
Then $\alpha\le\eta$, $\beta\le\eta$, and $\gamma\le\eta$.
The last inequality follows from $s\ge k/\eta$.
Since $\eta<1/100$, we have $r_0\le8\eta$ and $c_0\le11\eta$.
Using $\sqrt{OS}\le(O+S)/2$ in \eqref{eq:globalineq} gives
\[
 S\le\frac{1+r_0+\tfrac32c_0}{1-r_0-\tfrac12c_0}O
 \le(1+50\eta)O.
\]
The same inequality gives $S=0$ when $O=0$; division by $O$ is unnecessary.

For the final size-$2k$ diagonal blocks, choose
\begin{equation}\label{eq:leafwidth}
 p=2k+1+\lceil2k/\eta\rceil
\end{equation}
and query $A\Omega$ for an independent standard Gaussian
$\Omega\in\R^{n\times p}$. Subtract the now-known off-diagonal
approximation $H_{\rm off}\Omega$. On a diagonal block $I$, set
\[
 \widehat D_I=((A-H_{\rm off})\Omega)_I\,\Omega_I^\dagger.
\]
The contaminating columns outside $I$ are independent of $\Omega_I$.
The same Gaussian inverse moment calculation gives conditional total
diagonal error at most
$2k/(p-2k-1)$ times the off-diagonal error. Thus the final squared error obeys
\begin{equation}\label{eq:expectfinal}
 \E\norm{A-\widehat H}_\F^2
 \le(1+\eta)(1+50\eta)O\le(1+52\eta)O.
\end{equation}
Every output block meets its rank constraint, so $\widehat H\in\Hclass$.

\begin{theorem}[Explicit nonadaptive upper bound]\label{thm:upper}
With \eqref{eq:parameters}--\eqref{eq:leafwidth}, the preceding construction
uses at most
\begin{equation}\label{eq:budget}
 Q=2(L-1)^2(s+t)+p
\end{equation}
queries and satisfies \eqref{eq:guarantee}. Selecting the $n$-query exact
baseline whenever $Q\ge n$ gives
\[
 q_*\le\min\{n,Q\}=O\!\left(\min\{n,kL^2/\eps^2\}\right).
\]
\end{theorem}
\begin{proof}
The query count follows directly from the two orientations, buckets, sketch
widths, and last diagonal fit. If $O>0$, the excess
$Z=\norm{A-\widehat H}_\F^2-O$ is nonnegative by
Lemma~\ref{lem:separable}. Equation~\eqref{eq:expectfinal} and Markov's
inequality give
\[
 \Prob\{Z>(2\eps+\eps^2)O\}
 \le\frac{52\eta}{2\eps+\eps^2}\le0.01.
\]
If $O=0$, \eqref{eq:expectfinal} says the output error is zero almost surely.
Finally $s=O(k/\eta)$, $t=O(k/\eta^2)$, and $p=O(k/\eta)$,
with constants independent of $n,k,L,\eps$.
\end{proof}

\section{Joint fitting when the column spaces are supplied}\label{sec:jointfit}

The repeated coefficient-fitting cost across levels is avoidable when the
correct block column spaces are supplied. This conditional result does not
supply those spaces and therefore is not an upper bound for the general
problem by itself.

For each ordered sibling block $(I,J)$, fix an orthonormal matrix
$U_{I,J}\in\R^{|I|\times r_{I,J}}$ with $0\le r_{I,J}\le k$.
Let $\mathcal V(U)$ be the linear subspace of matrices whose block $(I,J)$
has columns in $\col(U_{I,J})$, with arbitrary diagonal leaves.
Then $\mathcal V(U)\subseteq\Hclass$.
For a column index $j$, concatenate the corresponding embedded sibling
bases along its root-to-leaf path and the coordinate basis of its leaf.
Call the resulting matrix $Z_j$. The row supports of its constituent pieces
are disjoint, so $Z_j$ is orthonormal and
\[
 r_j:=\#\text{columns}(Z_j)
   =k+\sum_{(I,J):j\in J}r_{I,J}\le k(L+1)=:r_{\max}^{\rm bound}.
\]
Put $r_{\max}=\max_jr_j$ and
$A_U[:,j]=Z_jZ_j^{\mathsf T}A[:,j]$.
This is the Frobenius projection of $A$ onto $\mathcal V(U)$.

Draw a fresh standard Gaussian $\Omega\in\R^{n\times q}$ and use $q$
transpose queries to obtain $Y=A^{\mathsf T}\Omega$. Define
\begin{equation}\label{eq:jointfit}
 B[:,j]=Z_j(\Omega^{\mathsf T}Z_j)^\dagger Y[j,:]^{\mathsf T}.
\end{equation}
The same $q$ responses are used for every column and every level.

\begin{proposition}[Known-space coefficient fit]\label{prop:jointfit}
If $q>r_{\max}+1$, the fit \eqref{eq:jointfit} belongs to $\mathcal V(U)$ and
satisfies the exact expectation identity
\begin{equation}\label{eq:jointidentity}
 \E\norm{A-B}_\F^2
 =\sum_{j=1}^n\left(1+\frac{r_j}{q-r_j-1}\right)
      \norm{(I-Z_jZ_j^{\mathsf T})A[:,j]}^2.
\end{equation}
In particular, taking
\[
 q=r_{\max}+1+\lceil50r_{\max}/\eps\rceil
\]
gives $\norm{A-B}_\F\le(1+\eps)\norm{A-A_U}_\F$ with probability at least
$0.99$, using $O(kL/\eps)$ queries.
\end{proposition}
\begin{proof}
Write $A[:,j]=Z_jc_j+e_j$ with $Z_j^{\mathsf T}e_j=0$.
For each fixed column, $G_j=\Omega^{\mathsf T}Z_j$ is standard Gaussian
and independent of $\Omega^{\mathsf T}e_j$, whose entries are independent
$N(0,\norm{e_j}^2)$. The least-squares coefficient error is
$G_j^\dagger\Omega^{\mathsf T}e_j$.
Lemma~\ref{lem:regression} gives its expected squared norm as
$r_j\norm{e_j}^2/(q-r_j-1)$. Orthogonality to $e_j$ proves the identity
for that column, and linearity of expectation proves \eqref{eq:jointidentity}.
The fits for different columns need not be independent.

The error in excess of $\norm{A-A_U}_\F^2$ is nonnegative and has mean at
most $(\eps/50)\norm{A-A_U}_\F^2$ for the stated width.
Markov's inequality bounds the failure probability by
$1/(50(2+\eps))<0.01$. If the projection error is zero, the expectation
identity gives exact recovery almost surely.
\end{proof}

If the spaces are obtained using earlier queries, the proposition remains
valid conditionally when $\Omega$ is independent of those earlier queries.
If the supplied bases contain the true leading $k$ left singular vectors of
each block, then $\norm{A-A_U}_\F=\OPT(A)$.
Neither observation yields those bases from the permitted oracle at no cost.
The accompanying test helper that forms bases from a full SVD is expressly
privileged input preparation, not part of a claimed query algorithm.

This result isolates one missing step more precisely: the $kL/\eps$ scale is
sufficient for fitting all coefficients once suitable rank-$k$ spaces are
available. The present work has no uniformly accurate procedure for learning
such spaces at the total cost needed to match the combined lower bound.

\section{A two-stage improvement with shared regression}\label{sec:twostage}

The coefficient accuracy needed during peeling is lower than the accuracy
needed in the final answer. Keeping the larger range spaces, using only a
constant-factor pilot fit, and then refitting all blocks simultaneously gives
a stronger upper bound than Theorem~\ref{thm:upper}.

\subsection{A constant-accuracy pilot and accurate retained ranges}

Let $h=L-1$ and choose
\begin{equation}\label{eq:twostageparameters}
 \eta=\eps/500,\qquad s=k+1+\lceil k/\eta\rceil,\qquad
 t=s+1+1000s,\qquad B=h.
\end{equation}
Run the peeling construction of Section~\ref{sec:algorithm} on the first
$h$ levels with these parameters, but do not perform its final diagonal fit.
Retain every range basis $Q_b$ as well as the rank-$k$ pilot block used for
subtraction at later levels. Write $r_b=\dim\col(Q_b)\le s$.
The two orientations at a level still use a frozen pilot from coarser levels.

In the notation of \eqref{eq:levelbound},
$\alpha\le\eta<10^{-3}$, $\beta=10^{-3}$, and
$\gamma\le\sqrt{\eta/1000}$. Hence $r_0<0.01$ and $c_0<0.02$.
The global inequality \eqref{eq:globalineq} gives, for the expected
nonterminal pilot error $S$,
\begin{equation}\label{eq:pilotconstant}
 S\le\frac{1+r_0+\tfrac32c_0}{1-r_0-\tfrac12c_0}O<2O,
 \qquad O=\OPT(A)^2.
\end{equation}
This is a constant-accuracy statement about the pilot, not the final output.

For each target block $C_b$, define the retained-range miss
\[
 M_b^{\rm miss}=\norm{(I-Q_bQ_b^{\mathsf T})[C_b]_k}_\F^2.
\]
The right-space estimate \eqref{eq:rangemiss}, conditioned on the past and
then averaged over hashes, implies
\begin{equation}\label{eq:rangesum}
 \E\sum_b M_b^{\rm miss}
 \le\alpha\left(O+\frac1h\sum_{\ell=1}^h\E T_{\ell-1}\right)
 \le\eta(O+S)\le3\eta O.
\end{equation}
Let $I_Q$ denote the squared Frobenius error of the best rank-$k$ approximation within each
retained column space, with exact size-$2k$ diagonal blocks. By
\eqref{eq:idealrange},
\begin{equation}\label{eq:idealshared}
 \E I_Q\le(1+3\eta)O.
\end{equation}

\subsection{Shared regression in larger spaces}

Now condition on the entire pilot. For each column $j$, form an orthonormal
matrix $Z_j$ by concatenating its size-$2k$ diagonal coordinate basis and
all retained $Q_b$ on sibling row intervals for which $j$ belongs to the
column interval. The row supports are disjoint. Thus
\[
 r_j=\dim\col(Z_j)\le2k+sh=:r_{\rm bd}.
\]
Let $A_Q[:,j]=Z_jZ_j^{\mathsf T}A[:,j]$, $e_j=A[:,j]-A_Q[:,j]$, and
$R_Q^2=\sum_j\norm{e_j}^2$. This larger linear model has no rank-$k$
restriction on the retained coefficients. In particular,
\begin{equation}\label{eq:projectionvsproper}
 R_Q^2\le I_Q.
\end{equation}

Draw a fresh standard Gaussian $\Omega\in\R^{n\times p}$ and obtain
$A^{\mathsf T}\Omega$ with $p$ transpose queries. Fit each column by
$Z_j(\Omega^{\mathsf T}Z_j)^\dagger\Omega^{\mathsf T}A[:,j]$.
On every nonterminal sibling block, truncate this fitted block to rank $k$.
Leave the fitted size-$2k$ diagonal blocks unchanged. Call the resulting
proper HODLR matrix $\widehat H$.

\begin{lemma}[Row isotropy survives shared fitting]\label{lem:sharedisotropy}
Fix a retained block $b=(I,J)$ with $r=r_b>0$, and write its fitted coefficient
matrix before truncation as $X_b+E_b$, where $X_b=Q_b^{\mathsf T}A_{I,J}$.
For $p>\max_jr_j+1$, there is a positive semidefinite matrix $\Sigma_b$ such
that
\[
 \E[(E_b)_{i\alpha}(E_b)_{j\beta}\mid\text{pilot}]
   =\delta_{ij}(\Sigma_b)_{\alpha\beta},
 \qquad
 \tr\Sigma_b=\sum_{a\in J}\frac{\norm{e_a}^2}{p-r_a-1}.
\]
The expectation is over the fresh shared Gaussian sketch. Errors for
different columns or blocks are not claimed to be independent.
\end{lemma}
\begin{proof}
Embed $Q_b$ into $\R^n$ on rows $I$. For every $a\in J$, $Z_a$ contains
this same basis. Its other columns have row support outside $I$, and $e_a$
is orthogonal to the embedded $Q_b$.
Therefore $G=\Omega^{\mathsf T}Q_b$ is independent of all the other design
components and all noise vectors $\Omega^{\mathsf T}e_a$ needed for these
columns, simultaneously for $a\in J$.

For any fixed orthogonal $O_r\in\R^{r\times r}$, replacing $G$ by $GO_r$
leaves that joint Gaussian distribution invariant. The least-squares
coefficient-error block is then replaced by $O_r^{\mathsf T}E_b$, because
the transformation acts on these design coordinates and not on the residual
noise. Thus $E_b$ has a left-orthogonally invariant distribution.
All second moments are finite by Lemma~\ref{lem:regression}.
Sign changes of one row make cross-row second moments zero, and row
permutations make within-row covariance matrices equal. This proves the
claimed form with a positive semidefinite $\Sigma_b$.

For a fixed column $a$, the entire design
$\Omega^{\mathsf T}Z_a$ is standard Gaussian and independent of
$\Omega^{\mathsf T}e_a$. Lemma~\ref{lem:regression} says that every
coefficient coordinate has variance $\norm{e_a}^2/(p-r_a-1)$.
Summing these diagonal entries over $a\in J$ proves the trace identity.
\end{proof}

\begin{proposition}[Conditional proper-fit error]\label{prop:sharedrisk}
If $p\ge r_{\rm bd}+1+d$ for $d>0$, then
\begin{equation}\label{eq:sharedrisk}
 \E[\norm{A-\widehat H}_\F^2\mid\text{pilot}]
 \le I_Q+10\sqrt{\frac{kh}{d}}\sqrt{O R_Q^2}
      +\frac{6sh+2k}{d}R_Q^2.
\end{equation}
\end{proposition}
\begin{proof}
Use Lemma~\ref{lem:sharedisotropy} in
Corollary~\ref{cor:isotropic} for each block, and add the orthogonal error
outside its retained space. The coefficient tail
$\tau_k(X_b)$ is no larger than $\tau_k(C_b)$.
Put $S_b=\tr\Sigma_b$. Every column belongs to exactly one ordered sibling
column interval at each of the $h$ nonterminal levels. Hence
\[
 \sum_b S_b\le\frac{h}{d}\sum_j\norm{e_j}^2=\frac{h}{d}R_Q^2.
\]
Cauchy--Schwarz gives
$\sum_b\tau_k(C_b)\sqrt{kS_b}\le\sqrt{khOR_Q^2/d}$, and
$\sum_b r_bS_b\le shR_Q^2/d$.
If $r_b<k$, use the nontruncated case of
Corollary~\ref{cor:isotropic}; its stated upper bound remains valid.
Zero-dimensional blocks contribute no regression noise.

For a column $j$, its $2k$ diagonal coefficients have total expected squared
error $2k\norm{e_j}^2/(p-r_j-1)$. Summing these errors adds at most
$2kR_Q^2/d$. All block errors add on disjoint entries, proving the result.
\end{proof}

\subsection{The improved query bound}

Take
\begin{equation}\label{eq:sharedwidth}
 d=\left\lceil\max\left\{\frac{100kh}{\eta^2},
                    \frac{6sh+2k}{\eta}\right\}\right\rceil,
 \qquad p=2k+sh+1+d.
\end{equation}
Combining \eqref{eq:sharedrisk}, Jensen's inequality,
\eqref{eq:projectionvsproper}, and \eqref{eq:idealshared} gives
\begin{align*}
 \E\norm{A-\widehat H}_\F^2
 &\le \left(1+3\eta+\eta\sqrt{1+3\eta}
                      +\eta(1+3\eta)\right)O\le(1+6\eta)O.
\end{align*}
If $O=0$, the expected error is zero. Otherwise the final output is in
$\Hclass$, so its excess over $O$ is nonnegative. Markov's inequality yields
failure probability at most
$6\eta/(2\eps+\eps^2)\le0.006<0.01$.

\begin{theorem}[Two-stage nonadaptive upper bound]\label{thm:twostage}
The parameters \eqref{eq:twostageparameters} and
\eqref{eq:sharedwidth} give an algorithm satisfying
\eqref{eq:guarantee} with at most
\begin{equation}\label{eq:twostagebudget}
 Q_2=2(L-1)^2(s+t)+p
\end{equation}
queries. Using exact reconstruction when $Q_2\ge n$ gives
\[
 q_*(n,k,\eps)\le\min\{n,Q_2\}
 =O\left(\min\left\{n,\frac{kL^2}{\eps}
                            +\frac{kL}{\eps^2}\right\}\right).
\]
All query vectors may be drawn in advance. In particular, the final shared
Gaussian sketch has a width bounded using $r_{\rm bd}$, not a width selected
from the observed pilot ranks.
\end{theorem}
\begin{proof}
The pilot uses at most $2hB(s+t)=2h^2(s+t)$ queries and the shared fit uses
$p$ transpose queries. Here $s,t=O(k/\eps)$ and
$p=O(kL/\eps^2)$. The probability bound was established above.
The pilot's query vectors are nonadaptive by the construction in
Section~\ref{sec:algorithm}; the shared sketch is independent of all of them.
No oracle products with individual blocks are treated as free operations.
\end{proof}

The extra factor compared with the combined lower bound is now only one
factor of depth: before taking the $n$-query cap, the expression in this
upper bound is exactly $L(kL/\eps+k/\eps^2)$. A matching bound without this
factor is still not proved.


\section{Information in adaptive two-sided Gaussian queries}\label{sec:adaptive}

The next lower-bound construction is nonsymmetric and uses independent
Gaussian entries rather than GOE. Both query types must therefore be included
in the information calculation.

\begin{lemma}[Two-sided transcript divergence]\label{lem:rectkl}
Let $A_0,S\in\R^{n\times n}$ be fixed with $\norm S_\op\le1$ and let $G$
have independent $N(0,1)$ entries. For any randomized adaptive algorithm using
at most $q$ queries, let $\mathbb P_S$ and $\mathbb P_0$ be its transcript laws
on $A_0+\delta S+G$ and $A_0+G$. Include the seed in the transcript. Then
\begin{equation}\label{eq:rectkl}
 \KL(\mathbb P_S\Vert\mathbb P_0)\le\delta^2q/2.
\end{equation}
\end{lemma}
\begin{proof}
Keep separate orthonormal spans $V$ of forward query vectors and $U$ of
transpose query vectors. A component already in the appropriate span has a
known image; remove it and normalize the remainder. Redundant queries
contribute no new divergence.

For a fixed value of $S$ and a fixed history, the observations reveal $GV$ and
$U^{\mathsf T}G$. The unobserved component is
$P_{U^\perp}GP_{V^\perp}$. It remains a centered isotropic Gaussian on this
rectangular subspace. This is proved inductively: conditioned on the past,
the next direction is fixed, and the newly exposed row or column of the
remaining Gaussian rectangle is independent of the smaller rectangle left
unexposed. The same argument works for measurable adaptive choices because
they impose no condition beyond the answers already in the history.

For a fresh unit forward vector $v\perp V$, project the output onto
$U^\perp$. The omitted $U$ component is already determined by the known
transpose answers. At a common history the two conditional means of the
new information differ by $\delta P_{U^\perp}Sv$, and the conditional
covariance is the identity on $U^\perp$. Its one-step divergence is at most
$\delta^2/2$. A fresh transpose query similarly has mean difference
$\delta P_{V^\perp}S^{\mathsf T}u$ and the same bound. Apply the chain rule
and then average over the seed. Early stopping can be padded or treated by
zero-contribution steps.
\end{proof}

\begin{remark}[Known diagonal scaling]
Let $D$ be any fixed invertible diagonal matrix and write $A=A'D$.
An algorithm for $A$ can be simulated using the same number of queries to
$A'$: a forward query uses $A'(Dv)$, and a transpose query returns
$D(A'^{\mathsf T}v)$. Thus the same transcript-divergence bound applies to algorithms run on
$A=(A_0+\delta S+G)D$. The approximation loss for this scaled family must
still be analyzed separately; diagonal scaling does not preserve the
Frobenius objective. Neither the size of $D$ nor the size of $A_0$ incurs query
cost.
\end{remark}

\section{A weighted hierarchical frame packing}\label{sec:packing}

Assume $L\ge6$ and put $h=L-4$. In this section the symbol $h$ is a
lower-bound construction parameter, not the upper algorithm's $L-1$.
Use only upper sibling blocks at levels $\ell=1,\ldots,h$. Their orders are
$m_\ell=n/2^\ell$, and the smallest is $16k$.
For a target block $(I_b,J_b)$, let $C_b$ be the first $k$ coordinate indices
of $J_b$. These column sets are disjoint across the entire hierarchy:
an ancestor's selected indices lie in the leftmost leaf of its right child,
and hence in no strictly descendant right child.

Choose an orthonormal $m_\ell\times k$ frame $U_b$ and put it in rows $I_b$
and columns $C_b$ of a square matrix $S$, setting every other entry to zero.
Construct these frames from fine levels to coarse levels. At a block with
row interval $I_b$ and order $m$, require its frame to be orthogonal to all
previously constructed frame columns supported inside $I_b$.
There are exactly
\[
 k(2^{h-\ell}-1)=m/16-k
\]
such descendant columns. Thus the available subspace has dimension
$15m/16+k$; even the last column of the new frame is chosen in dimension at
least $15m/16+1$.

All nonzero columns of $S$ are consequently orthonormal. In particular,
$\norm S_\op=1$ and $S\in\Hclass$. Set the fixed diagonal weights to
\begin{equation}\label{eq:weights}
 D_{jj}^2=\begin{cases}
 m_\ell/n,&j\in C_b\text{ for a level-}\ell\text{ target},\\
 1/n,&j\text{ is not selected}.
 \end{cases}
\end{equation}
The weights do not depend on the chosen frames. At every target level,
$\sum_b\norm{U_b D_{C_b}}_\F^2=k/2$, so
\begin{equation}\label{eq:weightednorm}
 \norm{SD}_\F^2=kh/2.
\end{equation}

\begin{lemma}[Hierarchical packing]\label{lem:depthpacking}
There is a finite family $\{S_\theta:\theta\in\mathcal T\}$ constructed above
such that, with $N=|\mathcal T|$,
\begin{equation}\label{eq:depthpack}
 \log N\ge nkh/128,\qquad
 \norm{(S_\theta-S_{\theta'})D}_\F^2>kh/128
 \quad(\theta\ne\theta').
\end{equation}
Here and below logarithms used in information bounds are natural.
\end{lemma}
\begin{proof}
Sample each frame uniformly in its allowed orthogonal complement,
conditionally on finer frames. We bound the measure of a small ball around
any fixed frame family $T$.

A newly sampled unit column $u$ lies on a sphere in dimension
$d\ge15m/16+1$. For any fixed unit vector $v$ in the ambient space,
$\norm{u-v}\le1/4$ implies $\ip uv\ge31/32$.
Projecting $v$ into the allowed subspace cannot enlarge this cap probability.
For a standard Gaussian representation of the sphere, put
$a=63/961<1/15$ and apply the exponential Markov bound with parameter $2$:
\[
 \Prob\{u_1\ge31/32\}
 \le (1-4a)^{-1/2}5^{-(d-1)/2}
 \le 2\,5^{-(d-1)/2}\le e^{-m/2}.
\]
The last inequality holds for $m\ge16$. This bound remains valid after any
previous columns and all finer frames have been fixed.

A level has $M=n/(2m)$ frames. If its weighted squared distance to $T$ is
at most $k/64$, its $Mk$ column distances have total square at most $Mk/32$.
At least half these columns must therefore have distance at most $1/4$.
A union bound over subsets of columns gives the uniform conditional bound
\[
 \Prob\{\text{this level has weighted distance}^2\le k/64
          \mid\text{finer frames}\}
 \le2^{Mk}e^{-mMk/4}\le e^{-nk/16}.
\]
If the total weighted squared distance is at most $kh/128$, at least half
of the $h$ levels have distance at most $k/64$.
For any specified collection of levels, integrate their conditional bounds
from coarse to fine (each is uniform in its finer conditioning).
A union bound over level subsets yields
\begin{equation}\label{eq:smallball}
 \Prob\{\norm{(S-T)D}_\F^2\le kh/128\}
 \le 2^h e^{-nkh/32}\le e^{-nkh/64}.
\end{equation}
The last inequality uses $nk\ge64$, valid for $L\ge6$.

A greedy construction can therefore choose at least
$\lfloor e^{nkh/64}\rfloor$ centers with pairwise distance greater than
$\sqrt{kh/128}$: before that many choices, the union of the forbidden balls
has measure less than one. Since $nkh/64\ge2$, the logarithm of this packing
size is at least $nkh/128$.
\end{proof}

\section{Anchors and the depth--accuracy lower bound}\label{sec:depthlower}

\subsection{A deterministic anchor inequality}

\begin{lemma}[Selected-column loss]\label{lem:anchor}
Write $C=[X\ Y]$ with $X\in\R^{m\times k}$ and
$\sigma_{\min}(X)\ge s_0>0$. Let $B=[B_1\ B_2]$ have rank at most $k$.
If $\norm{X-B_1}_\F\le K$ and $s_0\ge2K$, then
\begin{equation}\label{eq:anchor}
 \norm{C-B}_\F^2-\tau_k(C)^2
 \ge\norm{X-B_1}_\F^2-\frac{2K}{s_0}\norm Y_\F^2.
\end{equation}
\end{lemma}
\begin{proof}
The perturbation bound on the smallest singular value makes $B_1$ full
column rank. Thus every column of $B_2$ belongs to $\col(B_1)$.
Let $P_X,P_B$ be the projections onto $\col(X)$ and $\col(B_1)$.
The competitor $[X\ P_XY]$ has rank at most $k$, while the best possible
fit of $Y$ using $\col(B_1)$ is $P_BY$. It follows that the left side of
\eqref{eq:anchor} is at least
\[
 \norm{X-B_1}_\F^2+\tr\bigl(Y^{\mathsf T}(P_X-P_B)Y\bigr).
\]
For equal-dimensional subspaces,
$\norm{P_X-P_B}_\op=\norm{(I-P_B)P_X}_\op$.
Since $(I-P_B)X=(I-P_B)(X-B_1)$,
this is at most $K/s_0$, and in particular at most $2K/s_0$.
Substitution proves the displayed inequality.
\end{proof}

\subsection{The hard input distribution}

For the core argument assume
\begin{equation}\label{eq:depthcore}
 h=L-4\ge4,\qquad \eps h\ge1.
\end{equation}
Choose $S_\theta$ uniformly from Lemma~\ref{lem:depthpacking}.
Let $G$ be an independent standard iid Gaussian square matrix.
For each target choose a fixed orthonormal frame $U_b^0$ and construct $A_0$
by putting $M_0U_b^0$ in rows $I_b$ and selected columns $C_b$, with zeros
elsewhere. Set
\begin{equation}\label{eq:depthinput}
 A=(A_0+\delta S_\theta+G)D,
 \qquad \delta^2=2^{18}\eps n.
\end{equation}
The large number $M_0$ is fixed below, independently of $\theta,G$.
Since every clean target block has at most $k$ nonzero columns,
$(A_0+\delta S_\theta)D\in\Hclass$. Hence
\begin{equation}\label{eq:depthopt}
 \OPT(A)^2\le\norm{GD}_\F^2.
\end{equation}

There are $K_{\rm col}=k(2^h-1)=n/16-k$ selected columns. From
\eqref{eq:weights},
\[
 \E\norm{GD}_\F^2
 =n\tr D^2=n\left(\frac{kh}{2}+\frac{n-K_{\rm col}}n\right)\le nkh.
\]
On the selected entries of target blocks, the expected noise energy is
\begin{equation}\label{eq:targetnoise}
 \sum_b\E\norm{(GD)_{I_b,C_b}}_\F^2
 =\frac{k}{2}\sum_{\ell=1}^h m_\ell
 =\frac{kn}{2}(1-2^{-h})<kn/2.
\end{equation}
Markov's inequality and a union bound therefore give an event of probability
at least $0.8$ on which
\begin{equation}\label{eq:depthgood}
 \norm{GD}_\F^2\le N_0:=10nkh,
 \qquad \sum_b\norm{(GD)_{I_b,C_b}}_\F^2\le5nk.
\end{equation}

Put
\[
 K_0=\tfrac32\sqrt{N_0},\quad s_0=2K_0/\eps,
 \quad w_{\min}=\sqrt{16k/n},
 \quad M_0=\delta+\frac{s_0+\sqrt{N_0}}{w_{\min}}+1.
\]
On \eqref{eq:depthgood}, the selected columns $X_b=A_{I_b,C_b}$ have
\[
 \sigma_{\min}(X_b)
 \ge(M_0-\delta)\sqrt{m_\ell/n}-\sqrt{N_0}\ge s_0.
\]
The unselected columns $Y_b$ in each target block contain noise only.

\subsection{Decoding from a successful approximation}

Let an algorithm satisfying \eqref{eq:guarantee} return $B$.
On its success event and \eqref{eq:depthgood}, every selected-column error is
at most $K_0$, because the entire error is at most
$(1+\eps)\sqrt{N_0}<K_0$. Apply Lemma~\ref{lem:anchor} to each target block.
Its remainder coefficient $2K_0/s_0$ equals $\eps$.
By nonnegativity of all other local excesses,
\begin{align}\label{eq:selectederror}
 \sum_b\norm{X_b-B_{I_b,C_b}}_\F^2
 &\le\bigl(\norm{A-B}_\F^2-\OPT(A)^2\bigr)
       +\eps\sum_b\norm{Y_b}_\F^2\\
 &\le(2\eps+\eps^2)N_0+\eps N_0
 \le35\eps nkh.\nonumber
\end{align}
Subtract the fixed anchor on these selected entries and set all other entries
to zero. This gives an estimate $\widehat W$ of $\delta S_\theta D$ with
\begin{align}\label{eq:depthradius}
 \norm{\widehat W-\delta S_\theta D}_\F^2
 &\le2\sum_b\norm{X_b-B_{I_b,C_b}}_\F^2
       +2\sum_b\norm{(GD)_{I_b,C_b}}_\F^2\\
 &\le70\eps nkh+10nk\le80\eps nkh.\nonumber
\end{align}
The last inequality uses $\eps h\ge1$.
The packing's half-separation squared is larger than
\[
 \delta^2kh/512=512\eps nkh>80\eps nkh.
\]
Thus nearest-center decoding of $\widehat W$ identifies $\theta$ on that
joint event. The average decoding error is at most $0.21$: the noise event
fails with probability at most $0.2$, and the algorithm fails with probability
at most $0.01$ after averaging its pointwise guarantee over the hard inputs.

\subsection{The information contradiction}

Known scaling by $D$ does not change the query count. Lemma~\ref{lem:rectkl}
therefore gives, for the full transcript $\mathsf T$,
\[
 I(\theta;\mathsf T)\le\delta^2q/2.
\]
To see this directly, average the divergences to the common baseline
$A_0+G$. That average equals $I(\theta;\mathsf T)$ plus the nonnegative
divergence of the mixture transcript law to the baseline.

For a uniform label in a family of size $N$, every decoder obeys
\begin{equation}\label{eq:fanonew}
 p_e\ge1-\frac{I(\theta;\mathsf T)+\log2}{\log N}.
\end{equation}
This follows by conditioning on its error indicator:
$H(\theta\mid\mathsf T)\le\log2+p_e\log N$.
If
\[
 q\le\frac{kh}{1024\cdot2^{18}\eps},
\]
then $I(\theta;\mathsf T)\le nkh/2048\le(\log N)/16$.
In the core regime $\log N\ge nkh/128\ge8$, so \eqref{eq:fanonew} gives
$p_e>0.85$, contradicting $p_e\le0.21$.

\begin{proposition}[Depth--accuracy lower bound in the core]\label{prop:depthcore}
Under \eqref{eq:depthcore},
\[
 q_*(n,k,\eps)\ge c_{\rm d}\frac{kh}{\eps},
 \qquad c_{\rm d}=(1024\cdot2^{18})^{-1}.
\]
\end{proposition}

\section{Uniform extension and what is still missing}\label{sec:extension}

Appendix~\ref{app:accuracy} proves, using the earlier supplied argument,
\begin{equation}\label{eq:accuracyuniform}
 q_*(n,k,\eps)\ge c_*\min\{n,k/\eps^2\},
 \qquad c_*=(6400\cdot2^{48})^{-1}.
\end{equation}
We now remove the auxiliary restrictions from the new depth bound.

In the core, $h\ge L/2$ and
$kL/\eps\le kLh\le2kh^2\le n$ for $h\ge4$.
Thus Proposition~\ref{prop:depthcore} implies
$q_*\ge(c_{\rm d}/2)\min\{n,kL/\eps\}$.
Outside the core, either $L<8$, in which case $\eps L<4$, or $L\ge8$ and
$\eps h<1$, in which case $\eps L<1+4\eps<3$.
In both cases
\[
 kL/\eps\le4k/\eps^2.
\]
Equation~\eqref{eq:accuracyuniform} therefore yields the depth bound with
constant $c_*/4$. This is smaller than $c_{\rm d}/2$ and works in all cases.

For nonnegative $a,b$, one has
$\min(n,a)+\min(n,b)\ge\min(n,a+b)$.
Taking the maximum of the two established lower bounds and using half their
sum gives the left side of \eqref{eq:combined} with $c=c_*/8$.
The upper side follows from Theorem~\ref{thm:twostage}.
This completes the detailed argument for Theorem~\ref{thm:combined}.

\subsection{The roles of the two hard distributions}

The depth construction estimates many coefficients spread over scales. Its
weighted packing has entropy proportional to $nkh$, and large fixed anchors
convert small global excess into accurate estimation of selected columns.
The diagonal weighting keeps the information per whitened query controlled
while assigning comparable clean signal energy to each level.

The accuracy construction in the appendix is different. It places nearly
degenerate rank choices inside size-$2k$ sibling blocks. The relevant local
loss is linear in a small spectral gap, while transcript information is
quadratic in that gap. This produces $k/\eps^2$ rather than $k/\eps$.
The two constructions cannot simply be multiplied: the packing dimension,
noise normalization, and approximation-to-decoding inequalities all change
when scales are combined.

\subsection{The unresolved joint rate}

The expression
\[
 \min\{n,kL/\eps+k/\eps^2\}
\]
is a lower bound, not an established characterization. A matching upper
bound would need to use information across levels more efficiently than the
separate perforated regressions analyzed here. The two-stage algorithm reduces the final fitting cost by sharing a single
fresh regression across all levels. Its pilot still uses separate perforated
sketches at each level, giving the $kL^2/\eps$ term. The final rank-truncation
analysis has a $kL/\eps^2$ term because each column's residual enters $L-1$
blockwise noise sums. Together these leave only a single factor $L$ between
the proposed upper and lower expressions, but that factor is not a universal
constant.

\begin{proposition}[An obstruction to a uniform-embedding shortcut]\label{prop:nullspace}
Let $m=n/2$, $\Omega\in\R^{n\times s}$, and $\Psi\in\R^{n\times t}$,
where $s,t<m$. There is a nonzero rank-one $H\in\Hclass$ satisfying
$H\Omega=0$ and $\Psi^{\mathsf T}H=0$.
\end{proposition}
\begin{proof}
Let $I,J$ be the left and right root children. Choose unit vectors
$u\in\ker(\Psi_I^{\mathsf T})$ and
$v\in\ker(\Omega_J^{\mathsf T})$. These nullspaces are nontrivial by the
width assumptions. Put $uv^{\mathsf T}$ in block $(I,J)$ and zero elsewhere.
This matrix has Frobenius norm one, lies in $\Hclass$ for every $k\ge1$, and
both products vanish.
\end{proof}

This proposition rules out an isometry or lower norm bound holding
simultaneously for all HODLR differences under these two short sketches.
Its witness depends on the realized sketches. It is therefore \emph{not} a
lower bound against the pointwise randomized guarantee in \eqref{eq:guarantee},
and does not rule out a successful joint estimator with a different analysis.

Replacing the missing argument by an assertion that a joint least-squares
fit must be better would not be valid: the fitted subspaces depend on the
same sketch data, and rank truncation is nonlinear. Neither independence nor
a uniform embedding of all possible rank-constrained differences follows
from the calculations above. These are concrete gaps in that route, not
steps that have been assumed in Theorems~\ref{thm:upper} and~\ref{thm:twostage}.

\section{Reproducibility and proof-review priorities}\label{sec:checks}

The package includes a dense, query-counted NumPy reference implementation.
Its main routine uses only an oracle's dimension and vector-product method;
it does not read the oracle's stored matrix. Every column product, including
transpose products, is counted. The parameterized experimental routine may
use smaller widths, but those widths are explicitly distinguished from the
conservative theorem parameters.

The recorded checks use a fixed seed and include the following results.
\begin{center}\small
\begin{tabular}{@{}p{.56\linewidth}r@{}}
\toprule
Check & Instances/draws\\
\midrule
Deterministic rank-truncation inequality & 3,600\\
Gaussian inverse and row-isotropy moments & 12,000\\
Noisy Nystr\"om estimates (four noise scales) & 4,800\\
Hierarchical frame geometries & 5\\
Selected-column anchor inequality & 1,500\\
Fixed-query two-sided divergence identity & 120\\
Parameter and constant inequalities & 11,760\\
Oracle-level algorithm exercises & 13\\
\bottomrule
\end{tabular}
\end{center}
All checks in the recorded run passed. The largest observed hierarchical
orthogonality defect was below $1.6\times10^{-15}$, and the fixed-query
divergence identity discrepancies were below $7\times10^{-16}$.
The inverse-moment and row-second-moment Monte Carlo relative discrepancies
were approximately $0.0068$ and $0.0139$.

These tests do not certify a worst-case probabilistic theorem. In particular,
the fixed-query identity check does not test adaptivity, finite randomly
sampled frames do not establish packing entropy, and Monte Carlo averages
do not prove expectation inequalities. The adaptive posterior argument,
small-ball packing, all-parameter extension, and the exact-arithmetic
interpretation require mathematical review.

The two-stage refinement is implemented separately in \texttt{two\_stage.py}.
Its additional checks use 3,000 shared-regression second-moment trials and
400 conditional proper-fit trials. The second-moment discrepancies are below
1.8 recorded Monte Carlo standard errors. Twelve end-to-end cases check
proper output and oracle accounting; 840 parameter triples check the stated
constant inequalities. The actual two-stage path also recovers a diagonal
$1024\times1024$ example in 230 queries, with relative error below
$3.3\times10^{-16}$. These are experimental-width exact-instance checks,
not evidence of a uniform practical accuracy guarantee at those widths.

The separate known-space tests use 1,300 Gaussian trials and verify the
expectation identity in Proposition~\ref{prop:jointfit} within the recorded
Monte Carlo tolerances. They also recover a $128\times128$ known-space
instance in 20 queries with relative error below $1.5\times10^{-15}$.
Those tests supply spaces computed from explicit full matrices; this is not
an additional general-oracle guarantee.

A further audit checks 240 sketch-dependent nullspace witnesses and 90
adaptively selected two-sided Gaussian histories, comprising 1,320 query
steps. The largest discrepancy between the direct Gaussian posterior
projector and its rectangular-span formula is below $4.8\times10^{-14}$.
These finite-dimensional checks do not establish the conditional-independence
claim for all measurable adaptive algorithms. Four extreme-accuracy input
checks also exercise the $n$-query fallback through an opaque oracle, including
the smallest positive IEEE double-precision value. The earlier package's
verification suite was rerun without changing its source; its recorded checks
also pass.

For the algorithm, the small theorem-parameter examples use the exact
$n$-query fallback. Separately, smaller experimental widths exercise actual
peeling, including zero-tail inputs and nonzero perturbations. Their measured
error ratios are reported without claiming $0.99$ relative-error guarantees.
A diagonal $1024\times1024$ example is recovered in 260 queries to numerical
precision; this is only a query-count and exact-instance exercise.

Reproduce the checks from the package root with
\begin{lstlisting}
OPENBLAS_NUM_THREADS=1 OMP_NUM_THREADS=1 \
  python code/verify.py --output results/verification.json
OPENBLAS_NUM_THREADS=1 OMP_NUM_THREADS=1 \
  python code/verify_two_stage.py --output results/two_stage.json
\end{lstlisting}
The earlier archive is retained unchanged under \texttt{prior/}. The
appendix below incorporates its accuracy argument with only organizational
changes and a uniform-extension conclusion. Both the new proofs and that
earlier proof remain unreviewed.

\clearpage
\appendix
\section{The earlier accuracy lower bound}\label{app:accuracy}

This appendix supplies the $k/\eps^2$ argument used in the main theorem.
It retains the definitions and constants of the earlier note.

\subsection{Information in adaptive shifted-GOE queries}

For a symmetric matrix $G$, use the GOE normalization
\begin{equation}\label{eq:goe}
 G_{ij}=G_{ji}\sim N(0,1/n)\quad(i<j),\qquad
 G_{ii}\sim N(0,2/n),
\end{equation}
with independent upper-triangular entries. Its density on the vector space of
symmetric matrices is proportional to
$\exp(-n\norm{G}_{\F}^2/4)$. In particular, it is isotropic for the Frobenius
inner product on that vector space.

\begin{lemma}[Adaptive transcript divergence]\label{lem:kl}
Fix symmetric $A_0,S\in\R^{n\times n}$ with $\norm{S}_{\op}\le1$ and a real
$\delta$. Let $\mathbb P_S$ be the law of the transcript of any randomized
adaptive algorithm using at most $q\le n$ vector queries on
\[
 A=A_0+\delta S+G.
\]
Let $\mathbb P_0$ be its transcript law on $A_0+G$. Include the independent
random seed in the transcript. Then
\begin{equation}\label{eq:kl}
 \KL(\mathbb P_S\Vert\mathbb P_0)\le\frac{n\delta^2q}{2}.
\end{equation}
The statement permits queries on either side of $A$.
\end{lemma}

\begin{proof}
\emph{Reduction to orthonormal fresh directions.}
All matrices in question are symmetric, so choosing the transpose side gives
no additional type of measurement. Subtract the known product $A_0v$ from
each answer. If a new query vector has a component in the span of previous
queries, its image on that component is already known. Removing it and
normalizing the remaining vector recovers the same fresh information using
one call. Queries wholly in the old span give no information and can be
omitted. We may therefore work with an orthonormal sequence of fresh vectors,
chosen adaptively. An algorithm that stops early can be padded with
orthogonal queries. Data processing then bounds its original transcript by
the padded transcript. Fix the seed for the moment.

\emph{Gaussian exposure under adaptivity.}
After $t-1$ fresh directions, let $V$ be the matrix of those directions and
$\Pi=VV^{\mathsf T}$. For a fixed value of $S$, the observed answers determine
$Z=GV$. The symmetric component of $G$ determined by these observations is
\begin{equation}\label{eq:observedcomponent}
 G_{\rm obs}=ZV^{\mathsf T}+VZ^{\mathsf T}
             -V(V^{\mathsf T}Z)V^{\mathsf T}.
\end{equation}
The undetermined part is $(I-\Pi)G(I-\Pi)$. These are orthogonal subspaces of
the space of symmetric matrices for the Frobenius inner product. Conditional
on the history, the latter part remains a centered GOE compression on
$\operatorname{range}(I-\Pi)$, with the normalization $1/n$, not the inverse
of the remaining dimension.

For completeness, this claim also holds when $V$ is adaptive. Initially it
is just the orthogonal decomposition of an isotropic Gaussian. Given the
past, the next direction is fixed. In an orthonormal basis of the remaining
space beginning with that direction, the new Gaussian column and the
unexposed lower-right symmetric block are independent. Conditioning on the
column leaves that block unchanged. Induction proves the assertion. The
choice of the next vector is a function of already conditioned-on answers,
so it does not impose an additional constraint on the unexplored block.

\emph{One-step divergence.}
Let $v$ be the next unit vector, with $\Pi v=0$. The new information is the
component
\[
 y=(I-\Pi)(A-A_0)v.
\]
The component in $\operatorname{range}(\Pi)$ is already determined by the old
answers and symmetry:
$V^{\mathsf T}(A-A_0)v=((A-A_0)V)^{\mathsf T}v$.
At a fixed history, the conditional mean of $y$ under the shifted law is
$\delta(I-\Pi)Sv$, and under the baseline it is zero. In both cases its
conditional covariance, restricted to $\operatorname{range}(I-\Pi)$, is
\begin{equation}\label{eq:covariance}
 \Sigma_v=\frac1n\bigl((I-\Pi)+vv^{\mathsf T}\bigr).
\end{equation}
The variance is $2/n$ in direction $v$ and $1/n$ in the orthogonal remaining
directions. In particular, $\Sigma_v^{-1}\preceq nI$ on its support.

Two Gaussians with identical covariance and mean difference $m$ have
divergence $m^{\mathsf T}\Sigma_v^{-1}m/2$, as follows directly by integrating
the logarithm of their density ratio. The conditional divergence is therefore
at most
\[
 \frac{n\delta^2}{2}\norm{(I-\Pi)Sv}^2
 \le\frac{n\delta^2}{2}.
\]
Applying the chain rule for relative entropy to the successive fresh answers
and summing gives \eqref{eq:kl}. Integrating over the seed gives the same bound
for randomized algorithms, since its distribution is identical in both
experiments.
\end{proof}

\begin{remark}[Why exact arithmetic does not invalidate the bound]
This is a divergence calculation for the laws of exact real-valued answers.
No rounding, finite-precision hypothesis, bounded query coefficients, or
nonadaptive restriction is used. Normalizing a query loses no information
because its scale is known. The unexplored Gaussian compression is what
prevents arbitrary measurable postprocessing from revealing the hidden
parameter for free.
\end{remark}

\subsection{A local approximation-to-subspace inequality}

The next lemma is the deterministic part of the construction. It handles
arbitrary perturbations, nonsymmetric estimates, and estimates of rank
strictly smaller than $k$.

\begin{lemma}[Local loss]\label{lem:local}
Let $P$ be a rank-$k$ orthogonal projector on $\R^{2k}$, and let
$0<\delta\le1/2$. Put
\[
 C=(1-\delta)I+2\delta P.
\]
For any $W\in\R^{2k\times2k}$ and any matrix $D$ of rank at most $k$, choose
a rank-$k$ orthogonal projector $Q$ whose range contains the row space of $D$.
Then
\begin{align}\label{eq:local}
 &\norm{C+W-D}_{\F}^2-\tau_k(C+W)^2\notag\\
 &\hspace{1em}\ge
 \delta\norm{P-Q}_{\F}^2
 -\left(\frac{(1+\delta)^2}{\delta}+1\right)\norm{W}_{\F}^2.
\end{align}
\end{lemma}

\begin{proof}
Write $d=\norm{P-Q}_{\F}$ and $w=\norm{W}_{\F}$. Since $D=DQ$,
\[
 \norm{C+W-D}_{\F}^2\ge\norm{(C+W)(I-Q)}_{\F}^2.
\]
The rank-$k$ matrix $(1+\delta)P=CP$ is a valid competitor, so
\[
 \tau_k(C+W)^2\le\norm{C(I-P)+W}_{\F}^2.
\]
Subtracting these inequalities and expanding squared norms yields
\begin{align}\label{eq:localexpand}
 &\norm{C(I-Q)}_{\F}^2-\norm{C(I-P)}_{\F}^2
  +2\ip{W}{C(P-Q)}_{\F}-\norm{WQ}_{\F}^2.
\end{align}
Because $C^2=(1-\delta)^2I+4\delta P$, the clean difference is
\[
 \tr((P-Q)C^2)=4\delta(k-\tr(PQ))=2\delta d^2.
\]
The cross term is bounded below by $-2(1+\delta)wd$, and the last term is at
least $-w^2$. Thus \eqref{eq:localexpand} is bounded below by
\[
 2\delta d^2-2(1+\delta)wd-w^2.
\]
Use
$2(1+\delta)wd\le\delta d^2+(1+\delta)^2w^2/\delta$
to obtain \eqref{eq:local}.
\end{proof}

The projector $Q$ can be selected measurably from $D$: choose a deterministic
orthonormal basis of its row space, then extend it by a fixed coordinate-based
Gram--Schmidt rule until its dimension is $k$. No optimality, symmetry, or
uniqueness assumption on $D$ is needed.

\subsection{An explicit finite packing}

All logarithms in this section and in the information calculation are natural.
The metric on tuples of projectors is the square root of the sum of their
squared Frobenius distances.

\begin{lemma}[Projector tuples]\label{lem:packing}
For every pair of positive integers $k,M$, there is a finite collection
\[
 \bigl\{(P_1^\theta,\ldots,P_M^\theta):\theta=1,\ldots,N\bigr\}
\]
of rank-$k$ orthogonal projectors on $\R^{2k}$ such that
\begin{equation}\label{eq:packsize}
 \log N\ge\frac{Mk^2}{200}
\end{equation}
and, whenever $\theta\ne\theta'$,
\begin{equation}\label{eq:packdist}
 \sum_{j=1}^M\norm{P_j^\theta-P_j^{\theta'}}_{\F}^2
 \ge\frac{Mk}{1600}.
\end{equation}
\end{lemma}

\begin{proof}
\emph{Step 1: a single-block alphabet.}
Suppose first that $k\ge4$. Let $S$ range over sign matrices with independent
entries $\pm1/\sqrt{k}$. For fixed unit vectors $x,y$, the random variable
$x^{\mathsf T}Sy$ has moment-generating function at most
$\exp(t^2/(2k))$, because $\cosh u\le\exp(u^2/2)$ and
$\sum_{ij}x_i^2y_j^2=1$. Thus
\[
 \Prob\{|x^{\mathsf T}Sy|>a\}\le2\exp(-ka^2/2).
\]
A $1/4$-net of the unit sphere has size at most $9^k$: take a maximal separated
set and compare disjoint balls of radius $1/8$ with the enclosing ball of
radius $9/8$. For nets in both variables,
$\norm{S}_{\op}\le2\max_{x,y\text{ in nets}}|x^{\mathsf T}Sy|$.
Consequently,
\[
 \Prob\{\norm{S}_{\op}>10\}
 \le2\exp\bigl(-(25/2-2\log9)k\bigr)<1/2.
\]
At least $2^{k^2-1}$ sign matrices are therefore good.

Let $h(p)=-p\log p-(1-p)\log(1-p)$. A binary Hamming ball of radius
$\lfloor k^2/4\rfloor$ has size at most $\exp(h(1/4)k^2)$. To verify the usual
binomial bound, multiply each summand by
$p^i(1-p)^{d-i}$ with $d=k^2$ and $p=1/4$; for $i\le pd$ this weight is at
least $\exp(-dh(p))$, while the weighted sum is at most one.
Greedy deletion of such balls from the good matrices gives a set with mutual
Hamming distance greater than $k^2/4$ and log cardinality at least
\[
 (\log2-h(1/4))k^2-\log2\ge k^2/20,\qquad k\ge4.
\]

For each retained sign matrix set $T=S/20$, so $\norm{T}_{\op}\le1/2$, and let
$P_T$ project onto the graph of $T$, the column space of
$[I\;T^{\mathsf T}]^{\mathsf T}$. For two matrices $T,U$, an orthonormal basis
of the graph of $U$ is
$[I\;U^{\mathsf T}]^{\mathsf T}(I+U^{\mathsf T}U)^{-1/2}$; an orthonormal
basis of the orthogonal complement of the graph of $T$ is
$[-T\;I]^{\mathsf T}(I+TT^{\mathsf T})^{-1/2}$. Hence
\begin{align}\label{eq:graphidentity}
 \frac12\norm{P_T-P_U}_{\F}^2
 =\norm{(I+TT^{\mathsf T})^{-1/2}(U-T)
          (I+U^{\mathsf T}U)^{-1/2}}_{\F}^2.
\end{align}
Both inverse square roots have minimum singular value at least
$1/\sqrt{1+1/4}$. The Hamming separation also gives
$\norm{T-U}_{\F}^2>k/400$. It follows that
\[
 \norm{P_T-P_U}_{\F}^2
 \ge\frac{2}{(1+1/4)^2}\norm{T-U}_{\F}^2\ge k/400.
\]
For $k=1,2,3$, instead take the two complementary coordinate projectors.
Their squared distance is $2k$ and their log cardinality is $\log2\ge k^2/20$.
Thus for every $k$ there is a single-block alphabet of size $Q\ge2$ with
$\log Q\ge k^2/20$ and squared separation at least $k/400$.

\emph{Step 2: a product code.}
In the $Q$-ary cube of length $M$, a Hamming ball of radius $\lfloor M/4\rfloor$
has size at most
\[
 \exp\bigl(Mh(1/4)+(M/4)\log(Q-1)\bigr).
\]
Indeed, factor $(Q-1)^{M/4}$ out of the relevant binomial sum and apply the
binary bound just proved. Greedy packing then gives mutual Hamming distance
greater than $M/4$ and
\[
 \log N\ge M\bigl(\log Q-h(1/4)-\tfrac14\log(Q-1)\bigr)
 \ge\tfrac1{10}M\log Q\ge Mk^2/200.
\]
The middle inequality holds for $Q\ge2$: the function
$0.9\log Q-0.25\log(Q-1)$ is increasing there, and its value at $2$ is greater
than $h(1/4)$. Each differing coordinate contributes at least $k/400$ to
squared projector distance. This proves \eqref{eq:packdist}.
\end{proof}

\subsection{The hard distribution and the accuracy lower bound}

\subsubsection{Matrices indexed by the packing}

Set $M=n/(4k)=2^{L-2}$. Partition the coordinates into $M$ consecutive nodes
of size $4k$, and split each into two sets $I_j,J_j$ of size $2k$.
These are actual sibling blocks of the HODLR partition. Crucially, they are
not the unrestricted $k\times k$ leaf blocks.

Take the packing from Lemma~\ref{lem:packing}, and put
\begin{align}\label{eq:hardfamily}
 R_j^\theta&=2P_j^\theta-I_{2k},\notag\\
 A_0&=\operatorname{blockdiag}_{j=1}^M
       \begin{bmatrix}0&I_{2k}\\ I_{2k}&0\end{bmatrix},\notag\\
 S_\theta&=\operatorname{blockdiag}_{j=1}^M
       \begin{bmatrix}0&R_j^\theta\\ R_j^\theta&0\end{bmatrix},\qquad
 A=A_0+\delta S_\theta+G.
\end{align}
Each $R_j^\theta$ is a symmetric orthogonal matrix. Thus $S_\theta$ is symmetric
orthogonal too, and $\norm{S_\theta}_{\op}=1$.
The upper targeted block of the clean matrix is
\[
 C_j=(1-\delta)I_{2k}+2\delta P_j^\theta.
\]
Its singular values are $1+\delta$ and $1-\delta$, each repeated $k$ times.
The lower targeted block is its transpose. All other clean blocks are zero.
By Lemma~\ref{lem:separable},
\begin{equation}\label{eq:cleanopt}
 \OPT(A_0+\delta S_\theta)^2
 =2Mk(1-\delta)^2=\frac n2(1-\delta)^2.
\end{equation}
The packing parameters become
\begin{equation}\label{eq:packingn}
 \log N\ge\frac{nk}{800},\qquad
 \sum_j\norm{P_j^\theta-P_j^{\theta'}}_{\F}^2\ge\frac n{6400}.
\end{equation}

\subsubsection{A good-noise event}

Define
\[
 Z=\norm{G}_{\F}^2,\qquad
 T=\sum_{j=1}^M\norm{G_{I_j,J_j}}_{\F}^2.
\]
The GOE normalization gives
\[
 \E Z=n+1\le\tfrac54n,\qquad \E T=k.
\]
For the second identity, there are $M(2k)^2=nk$ entries in the targeted
\emph{upper} blocks, each with variance $1/n$. These entries are disjoint.
Markov's inequality and a union bound imply that the event
\begin{equation}\label{eq:event}
 \mathcal E=\{Z\le10n,\ T\le10k\}
\end{equation}
has probability at least
\begin{equation}\label{eq:eventprob}
 \Prob(\mathcal E)\ge1-\frac18-\frac1{10}=\frac{31}{40}.
\end{equation}
On this event, use the best clean HODLR approximation as a competitor for the
noisy matrix and apply $\norm{X+Y}_{\F}^2\le2\norm X_{\F}^2+2\norm Y_{\F}^2$:
\begin{equation}\label{eq:noisyopt}
 \OPT(A)^2\le n(1-\delta)^2+2Z\le21n.
\end{equation}
No independence between $Z$ and $T$ is needed.

\subsubsection{Decoding the hidden projectors}

Suppose an approximation algorithm succeeds on an input from
\eqref{eq:hardfamily}. Choose $Q_j$ to be a rank-$k$ projector containing the
row space of $B_{I_j,J_j}$, using the measurable rule following
Lemma~\ref{lem:local}. Let
\[
 r^2=\sum_{j=1}^M\norm{P_j^\theta-Q_j}_{\F}^2.
\]
Since $\eps<1/2$, success and \eqref{eq:noisyopt} imply, on $\mathcal E$,
\begin{align}\label{eq:globalexcess}
 \norm{A-B}_{\F}^2-\OPT(A)^2
 &\le (2\eps+\eps^2)\OPT(A)^2\notag\\
 &\le\frac{105}{2}\eps n.
\end{align}
Sum Lemma~\ref{lem:local} over the targeted upper blocks. Their excesses are
nonnegative and sum to at most the global excess by
Lemma~\ref{lem:separable}. Also,
\[
 \frac{(1+\delta)^2}{\delta}+1
 =\frac{1+3\delta+\delta^2}{\delta}\le\frac3\delta,
 \qquad 0<\delta\le1/2.
\]
It follows that
\begin{equation}\label{eq:radius}
 r^2\le\frac{105\eps n}{2\delta}+\frac{30k}{\delta^2}.
\end{equation}

Now set
\begin{equation}\label{eq:constants}
 D=2^{24},\qquad \eps_0=\frac1{2D},\qquad \delta=D\eps,
\end{equation}
and first restrict to
\begin{equation}\label{eq:coregime}
 0<\eps\le\eps_0,\qquad n\ge k/\eps^2.
\end{equation}
Then $0<\delta\le1/2$ and
\begin{equation}\label{eq:smallradius}
 \frac{r^2}{n}\le\frac{105}{2D}+\frac{30}{D^2}
 <\frac1{102400}<\frac1{25600}.
\end{equation}
The smallest squared separation in \eqref{eq:packingn} is at least $n/6400$.
Thus $r$ is less than half the minimum separation. Nearest-neighbor decoding
of the tuple $(Q_1,\ldots,Q_M)$ in the finite packing recovers $\theta$
uniquely. This is a measurable operation, and no query is needed to do it.

Let $\theta$ be uniform on the packing, independently of $G$ and the algorithm's
seed. The fixed-input guarantee \eqref{eq:guarantee} implies average success
at least $0.99$ under this distribution. The success event and $\mathcal E$
therefore intersect with probability at least
\begin{equation}\label{eq:decodeprob}
 \frac{31}{40}-\frac1{100}=\frac{153}{200}.
\end{equation}
Hence the decoding error probability is at most $47/200$.

\subsubsection{Information cannot support that decoding probability}

Let $\mathcal T$ denote the transcript including the seed. By
Lemma~\ref{lem:kl}, and using the baseline transcript law generated by $A_0+G$,
\begin{equation}\label{eq:mutualinfo}
 I(\theta;\mathcal T)
 \le\frac1N\sum_{\theta=1}^N
       \KL(\mathbb P_\theta\Vert\mathbb P_0)
 \le\frac{n\delta^2q}{2}.
\end{equation}
The first inequality follows from the identity
\[
 \frac1N\sum_\theta\KL(\mathbb P_\theta\Vert\mathbb P_0)
 =I(\theta;\mathcal T)+\KL(\overline{\mathbb P}\Vert\mathbb P_0),
\]
where $\overline{\mathbb P}$ is the mixture transcript law.

We recall the needed entropy bound directly. If $\widehat\theta$ is any
transcript-based decoder and $p_e=\Prob(\widehat\theta\ne\theta)$, conditioning
also on the error indicator gives
\[
 H(\theta\mid\mathcal T)\le\log2+p_e\log N.
\]
Since $H(\theta)=\log N$, this yields Fano's inequality in the form
\begin{equation}\label{eq:fano}
 p_e\ge1-\frac{I(\theta;\mathcal T)+\log2}{\log N}.
\end{equation}

Assume, for a contradiction, that
\begin{equation}\label{eq:qthreshold}
 q\le\frac{k}{6400\delta^2}.
\end{equation}
This bound is less than $n$ under \eqref{eq:coregime}, so the hypothesis
$q\le n$ in Lemma~\ref{lem:kl} applies. Equations \eqref{eq:mutualinfo} and
\eqref{eq:packingn} give $I(\theta;\mathcal T)\le(\log N)/16$.
Moreover, \eqref{eq:coregime} implies
$nk\ge4D^2$, so $\log N\ge nk/800\ge16\log2$.
Equation \eqref{eq:fano} therefore gives $p_e\ge7/8$, contradicting
$p_e\le47/200$ from \eqref{eq:decodeprob}. We have proved the following.

\begin{proposition}[Accuracy lower bound in the core regime]\label{prop:core}
If $\eps\le\eps_0$ and $n\ge k/\eps^2$, then
\begin{equation}\label{eq:corelower}
 q_*(n,k,\eps)\ge c_*\frac{k}{\eps^2},\qquad c_*=(6400D^2)^{-1}.
\end{equation}
\end{proposition}


\subsection{Extending the accuracy lower bound to all parameters}

Recall $D=2^{24}$, $\eps_0=1/(2D)$, and $c_*=(6400D^2)^{-1}$.
The function $q_*(n,k,\eps)$ is nonincreasing in $\eps$.
For $\eps\ge\eps_0$,
\[
 c_*\min\{n,k/\eps^2\}\le(c_*/\eps_0^2)k=k/1600,
\]
so Proposition~\ref{prop:exactlower} suffices.
For $\eps<\eps_0$ and $n\ge k/\eps^2$, apply Proposition~\ref{prop:core}.
Otherwise put $\zeta=\sqrt{k/n}>\eps$. If $\zeta\le\eps_0$, monotonicity
and the core bound at accuracy $\zeta$ give $q_*\ge c_*k/\zeta^2=c_*n$.
If $\zeta>\eps_0$, then $n<k/\eps_0^2$, and the exact-recovery obstruction
again gives $q_*\ge c_*n$. This proves
\[
 q_*(n,k,\eps)\ge c_*\min\{n,k/\eps^2\}
\]
for every parameter triple in RE-03. Together with the $n$-query baseline,
it also proves Corollary~\ref{cor:full}.


\begin{thebibliography}{9}
\bibitem{re03}
OpenProblemsInNLA. \emph{RE-03: Optimal matvec query complexity of HODLR
approximation}. Problem entry, audit dated September 10, 2026.
\url{https://github.com/ajt60gaibb/OpenProblemsInNLA/tree/main/randomized-and-low-rank-approximation/RE-03}.
Public statement checked September 13, 2026.

\bibitem{chen}
Tyler Chen, Feyza Duman Keles, Diana Halikias, Cameron Musco, Christopher Musco,
and David Persson.
\emph{Near-optimal hierarchical matrix approximation from matrix-vector
products}. SODA 2025, pp.~2656--2692; arXiv:2407.04686v2, October 24, 2024.
\url{https://arxiv.org/html/2407.04686v2}.
Theorems 1.1--1.2, Sections 3--6 and 8.

\bibitem{prior}
\emph{Adaptive-query lower bounds for RE-03}.
Earlier supplied unreviewed research note, 13 pages, in
\texttt{prior/RE03\_partial\_results.zip} and
\texttt{prior/RE03\_partial\_results.pdf}.
The appendix reproduces its relevant lower-bound proof rather than treating
it as an independently established external result.

\bibitem{halko}
Nathan Halko, Per-Gunnar Martinsson, and Joel A. Tropp.
\emph{Finding structure with randomness: Probabilistic algorithms for
constructing approximate matrix decompositions}.
SIAM Review 53(2), 217--288, 2011; arXiv:0909.4061.
\url{https://arxiv.org/abs/0909.4061}.
Background on Gaussian range finding; the moment calculations used here are
proved directly in the text.
\end{thebibliography}
\end{document}
